%!TEX program = xelatex
%!TEX encoding = UTF-8

\documentclass{beamer}
\usetheme{Madrid}
\usecolortheme{beaver}
% \mode<presentation>
\usepackage[singlelinecheck=false]{caption}
\usepackage{textpos}
\usepackage{amssymb}
\usepackage{amsmath,amsfonts,amsthm,amstext}
\usepackage[justification=centering]{caption}
\usepackage[brazil]{babel}
% \usepackage[latin1]{inputenc}
\usepackage{graphicx}
\graphicspath{{../../figuras_geogebra/}{../../Pictures/}}
\usepackage{enumitem}
\usepackage{multicol}
\usepackage{answers}
\usepackage{qrcode}

\usepackage{pifont}
\newcommand{\cmark}{\ding{228}}%

% requires version 0.3 of the package
\usepackage[customcolors,beamer]{hf-tikz}

\tikzset{style green/.style={
        set fill color=green!50!lime!60,
        set border color=white,
    },
    style cyan/.style={
        set fill color=cyan!90!blue!60,
        set border color=white,
    },
    style orange/.style={
        set fill color=orange!80!red!60,
        set border color=white,
    },
    hor/.style={
        above left offset={-0.15,0.31},
        below right offset={0.15,-0.125},
        #1
    },
    ver/.style={
        above left offset={-0.1,0.3},
        below right offset={0.15,-0.15},
        #1
    }
}

\usepackage{tikz,ifthen,tikz-3dplot}
\usepackage{siunitx}
\usetikzlibrary{circuits.ee.IEC}
\usetikzlibrary{lindenmayersystems}
\usetikzlibrary[shadings]
\usepackage{pgfplots}
\pgfplotsset{compat=1.15}
\usepackage{mathrsfs}
\usetikzlibrary{arrows}
\usepackage{tcolorbox}
\newcounter{saveenumi}
\newcommand{\seti}{\setcounter{saveenumi}{\value{enumi}}}
\newcommand{\conti}{\setcounter{enumi}{\value{saveenumi}}}

\resetcounteronoverlays{saveenumi}


\renewcommand{\qedsymbol}{$\blacksquare$}

%\setbeamertemplate{itemize items}[ball]

%\setbeamertemplate{itemize items}[default]
%-------------amatrix
% Augmented matrix.  Usage (note the argument does not count the aug col):
% \begin{amatrix}{2}
    %   1  2  3 \\  4  5  6
    % \end{amatrix}
\newenvironment{amatrix}[1]{%
    \left[\begin{array}{@{}*{#1}{r}|r@{}}
    }{%
    \end{array}\right]
}

\newtheorem{definicao}{Defini\c{c}\~ao}[section]
\newtheorem{definicoes}{Defini\c{c}\~oes}[section]
\newtheorem{exemplo}{Exemplo}[section]
\newtheorem{exemplos}{Exemplos}[section]
\newtheorem{exercicio}{Exerc{\'\i}cio}
\newtheorem{observacao}{Observa{\c c}{\~a}o:}[section]
\newtheorem{observacoes}{Observa{\c c}{\~o}es:}[section]
\newtheorem*{solucao}{Solu{\c c}{\~a}o:}
\newtheorem{proposicao}{Proposi\c{c}\~ao}
\newtheorem{lema}{Lema}
\newtheorem{teorema}{Teorema}
\newtheorem{corolario}{Corol\'ario}
\newenvironment{prova}[1][Prova]{\noindent\textbf{#1:} }{\qedsymbol}%{\ \rule{0.5em}{0.5em}}
\def\ano{2023}
\def\semestre{2}
\def\turma{11}
\newcommand{\nsub}{\varsubsetneq}
\newcommand{\vaz}{\emptyset}
\newcommand{\dom}{{\rm dom\,}}
\newcommand{\im}{{\rm Im\,}}
\newcommand{\sub}{\subseteq}
\newcommand{\n}{\mathbb{N}}
\newcommand{\z}{\mathbb{Z}}
\newcommand{\rac}{\mathbb{Q}}
\newcommand{\real}{\mathbb{R}}
\newcommand{\complex}{\mathbb{C}}
\newcommand{\cp}[1]{\mathbb{#1}}
\newcommand{\ch}{\mbox{\textrm{car\,}}\nobreak}
\newcommand{\p}[1]{\mbox{\textrm{posto(#1)}}}
\newcommand{\vesp}[1]{\vspace{ #1  cm}}
\newcommand{\compcent}[1]{\vcenter{\hbox{$#1\circ$}}}
\newcommand{\comp}{\mathbin{\mathchoice
{\compcent\scriptstyle}{\compcent\scriptstyle}
{\compcent\scriptscriptstyle}{\compcent\scriptscriptstyle}}}
\newcommand\T{\rule{0pt}{2.6ex}}

\title{Matriz}
\author[\autor]{\autor}
\institute[\instituto]{\instituto}
\date{}

\begin{document}
    \begin{frame}
        \maketitle
    \end{frame}

    \logo{\includegraphics[scale=.1]{logo-MAT.png}\vspace*{8.5cm}}

    \begin{frame}
        Vamos trabalhar também com o conjunto de matrizes com entradas em $\cp{K}$:\pause
        \[
            M_{p \times q}(\cp{K}) =
            \left\{
                \begin{bmatrix}
                    a_{11} & a_{12} & \cdots & a_{1q}\\
                    a_{21} & a_{22} & \cdots & a_{2q}\\
                    \vdots & & & \vdots\\
                    a_{p1} & a_{p2} & \cdots & a_{pq}
                \end{bmatrix}
                \mid a_{ij} \in \cp{K},\ 1 \le i \le p,\ 1 \le j \le q
            \right\}
        \]
    \end{frame}

    \begin{frame}
        Em $M_{p\times q}(\cp{K})$ vamos admitir definidas as operações de soma de matrizes  e multiplicação de matriz por escalar.

        Vamos também admitir as seguintes propriedades como verdadeiras:
        \begin{proposicao}
            Dadas matrizes $A$, $B$ e $C \in M_{p\times q}(\cp{K})$  e escalares $\alpha$, $\lambda \in \cp{K}$ vale:
            \begin{enumerate}[label={\roman*})]
                \item $(A + B) + C = A + (B + C)$
                \item $A + B = B + A$
                \item $0_{p \times q} + A = A  = A + 0_{p \times q}$,  onde $0_{p \times q}$ é matriz cujas entradas são todas zero.
                \seti
            \end{enumerate}
        \end{proposicao}
    \end{frame}

    \begin{frame}
        \begin{proposicao}
            \begin{enumerate}[label={\roman*})]
                \conti
                \item $A - A = 0_{p \times q}$
                \item $(\alpha\lambda) A = \alpha(\lambda A)  = \lambda(\alpha A)$
                \item $(\alpha + \lambda)A =  \alpha A + \lambda A$
                \seti
            \end{enumerate}
        \end{proposicao}
    \end{frame}

    \begin{frame}
        \begin{proposicao}
            \begin{enumerate}[label={\roman*})]
                \conti
                \item $\alpha(A + B) =  \alpha A + \alpha B$
                \item $1A = A$ e $(-1)A = -A$
                \item  $0A = 0_{p \times q}$
            \end{enumerate}
        \end{proposicao}
    \end{frame}

    \begin{frame}
        Para o produto de matrizes vamos admitir as seguintes propriedades como verdadeiras:
        \begin{proposicao}
            Sejam $A$, $B$ e $C$  matrizes de ordens tais que as operações a seguir estejam definidas. Então:
                \begin{enumerate}[label={\roman*})]
                    \item $(A\cdot B)\cdot C = A\cdot(B \cdot C)$
                    \item $(A + B)\cdot C = (A\cdot C) + (B\cdot C)$
                    \item $C\cdot(A + B) = (C\cdot A) + (C\cdot B)$
                    \seti
                \end{enumerate}
        \end{proposicao}
    \end{frame}

    \begin{frame}
        \begin{proposicao}
            \begin{enumerate}[label={\roman*})]
                \conti
                \item $\alpha(A\cdot B) =  (\alpha A)\cdot B  = A \cdot(\alpha B)$
                \item Se $A$ é uma matriz quadrada de ordem $n$,  então \[A\cdot I_n = A = I_n\cdot A,\]  onde $I_n$ é a matriz quadrada de ordem $n$  onde todas as entradas da diagonal principal são 1  e as demais são zero.
                \item Se $A$ é uma matriz quadrada de ordem $n$,  então \[A\cdot0_n = 0_n  = 0_n\cdot A.\]
            \end{enumerate}
        \end{proposicao}
    \end{frame}
\end{document}
